% ==============================================================================
% Plantilla Institucional de Trabajo Terminal — UPIIZ IPN
% Archivo: capitulos/protocolo/09-viabilidad.tex
% ==============================================================================
% ESTRUCTURA NORMATIVA DE PROTOCOLO: CAPÍTULO 9 - VIABILIDAD DEL PROYECTO
% ==============================================================================
% LÍMITE INSTITUCIONAL: Máximo 5 páginas.
%
% Requisitos normativos (Anexo 1):
% 9.1 Recursos humanos:
%     - Alumnos participantes y datos de contacto.
%     - Asesores en orden jerárquico (máximo 3 asesores).
% 9.2 Equipo e instalaciones necesarias:
%     - Equipos, dispositivos, herramientas y laboratorios requeridos.
%     - Actividad asociada a cada recurso.
%     - Disponibilidad actual, compra o arrendamiento según corresponda.
% 9.3 Costo estimado y fuente de financiamiento:
%     - Desglose por etapas.
%     - Costo mínimo aproximado y costo máximo aproximado.
%     - Fuente de financiamiento y disponibilidad actual o forma prevista de obtenerlo.
% 9.4 Cronograma de actividades de Trabajo Terminal I y II:
%     - Calendario semanal.
%     - Actividades por alumno.
%     - Gráfica de Gantt (figuras externas mediante \includegraphics o tablas).
% - Bloque de firmas de conformidad del cronograma de actividades (Alumnos y Asesores).
% ==============================================================================

\chapter{Viabilidad del proyecto}
\label{cap:proto-viabilidad}

% [INSTRUCCIÓN PARA EL ALUMNO]:
% Desarrolle este capítulo en un máximo de 5 páginas cumpliendo puntualmente los cuatro apartados:

% ------------------------------------------------------------------------------
\section{Recursos humanos}
\label{sec:proto-recursos-humanos}

\subsection{Alumnos proponentes}
% Indique nombre completo, número de boleta y datos de contacto de cada alumno:

\begin{itemize}
    \item \textbf{Alumno(a) 1:} \alumnoANombre
    \begin{itemize}
        \item Boleta: \alumnoABoleta
        \item Correo institucional: [correo@alumno.ipn.mx]
        \item Teléfono de contacto: [Teléfono de contacto]
    \end{itemize}
    \ifcsvoid{alumnoBNombre}{}{%
        \item \textbf{Alumno(a) 2:} \alumnoBNombre
        \begin{itemize}
            \item Boleta: \alumnoBBoleta
            \item Correo institucional: [correo@alumno.ipn.mx]
            \item Teléfono de contacto: [Teléfono de contacto]
        \end{itemize}
    }%
    \ifcsvoid{alumnoCNombre}{}{%
        \item \textbf{Alumno(a) 3:} \alumnoCNombre
        \begin{itemize}
            \item Boleta: \alumnoCBoleta
            \item Correo institucional: [correo@alumno.ipn.mx]
            \item Teléfono de contacto: [Teléfono de contacto]
        \end{itemize}
    }%
\end{itemize}

\subsection{Asesores del proyecto (orden jerárquico, máximo 3)}
% Liste a los asesores en orden de prelación jerárquica conforme al Anexo 1:

\begin{enumerate}
    \item \textbf{Primer Asesor:} \asesorANombre
    \begin{itemize}
        \item Correo institucional: [correo@ipn.mx]
    \end{itemize}
    \ifcsvoid{asesorBNombre}{}{%
        \item \textbf{Segundo Asesor:} \asesorBNombre
        \begin{itemize}
            \item Correo institucional: [correo@ipn.mx]
        \end{itemize}
    }%
    \ifcsvoid{asesorCNombre}{}{%
        \item \textbf{Tercer Asesor:} \asesorCNombre
        \begin{itemize}
            \item Correo institucional: [correo@ipn.mx]
        \end{itemize}
    }%
\end{enumerate}

% ------------------------------------------------------------------------------
\section{Equipo e instalaciones necesarias}
\label{sec:proto-instalaciones}

% [INSTRUCCIÓN PARA EL ALUMNO]:
% Especifique en la siguiente tabla los recursos necesarios, la actividad en la que
% intervienen y su condición (disponible en UPIIZ, compra requerida o arrendamiento):

\begin{table}[htbp]
    \centering
    \caption{Relación de equipo, dispositivos, herramientas e instalaciones requeridas.}
    \label{tab:proto-recursos-materiales}
    \begin{tabularx}{\textwidth}{p{4.2cm} p{3.0cm} X c}
        \toprule
        \textbf{Equipo / Instalación} & \textbf{Tipo de recurso} & \textbf{Actividad asociada} & \textbf{Condición} \\
        \midrule
        {[Nombre del laboratorio]} & Instalación & [Pruebas o ensamble del sistema] & Disponible en UPIIZ \\
        {[Equipo o maquinaria]} & Equipo principal & [Maquinado o caracterización] & Disponible en UPIIZ \\
        {[Dispositivo / Sensor]} & Componente & [Medición de variables del proceso] & Compra requerida \\
        {[Herramienta específica]} & Herramienta & [Ensamble o ajuste mecánico] & Compra / Arrendamiento \\
        \bottomrule
    \end{tabularx}
\end{table}

% ------------------------------------------------------------------------------
\section{Costo estimado y fuente de financiamiento}
\label{sec:proto-costos-financiamiento}

% [INSTRUCCIÓN PARA EL ALUMNO]:
% Presente el desglose presupuestal estimado por etapas, especificando el costo
% mínimo aproximado y el costo máximo aproximado:

\begin{table}[htbp]
    \centering
    \caption{Presupuesto estimado por etapas del proyecto.}
    \label{tab:proto-presupuesto-etapas}
    \begin{tabularx}{\textwidth}{X c c}
        \toprule
        \textbf{Etapa del Proyecto} & \textbf{Costo mínimo aprox. (MXN)} & \textbf{Costo máximo aprox. (MXN)} \\
        \midrule
        Etapa 1: Componentes y materiales de diseño & [\$ 0.00] & [\$ 0.00] \\
        Etapa 2: Fabricación, maquinado e insumos & [\$ 0.00] & [\$ 0.00] \\
        Etapa 3: Integración de circuitos y sensores & [\$ 0.00] & [\$ 0.00] \\
        Etapa 4: Pruebas y validación experimental & [\$ 0.00] & [\$ 0.00] \\
        \midrule
        \textbf{Costo total estimado del proyecto:} & \textbf{[\$ 0.00]} & \textbf{[\$ 0.00]} \\
        \bottomrule
    \end{tabularx}
\end{table}

\noindent\textbf{Fuente de financiamiento y disponibilidad:}\\[0.3em]
[Especifique la fuente de financiamiento prevista (recursos propios, proyectos de investigación institucional SIP-IPN, apoyos de vinculación o fondos externos), indicando si el recurso se encuentra actualmente disponible o la forma prevista de obtenerlo].

% ------------------------------------------------------------------------------
\section{Cronograma de actividades de Trabajo Terminal I y II}
\label{sec:proto-cronograma-actividades}

% [INSTRUCCIÓN PARA EL ALUMNO]:
% El cronograma debe contemplar un calendario semanal detallando las actividades
% específicas asignadas a cada alumno para TT I y TT II.
% Puede incluirse mediante tablas estructuradas o figuras externas de gráficas de Gantt
% generadas con herramientas de gestión de proyectos mediante \includegraphics.

\subsection{Cronograma semanal de actividades de Trabajo Terminal I}

[Describa la planificación semanal de TT I, indicando la distribución de tareas entre los alumnos del equipo].

\begin{figure}[htbp]
    \centering
    % Inserte aquí su gráfica de Gantt exportada en PDF o PNG:
    % \includegraphics[width=0.95\textwidth]{figuras/proyecto/gantt-semanal-tt1.pdf}
    \fbox{\parbox[c][5cm][c]{0.9\textwidth}{\centering\textbf{[Espacio para Gráfica de Gantt Semanal de Trabajo Terminal I]}\\[0.5em]
    \small Inserte aquí el cronograma gráfico generado externamente mediante:\\
    \texttt{\textbackslash includegraphics[width=0.9\textbackslash textwidth]\{figuras/proyecto/gantt-semanal-tt1.pdf\}}}}
    \caption{Gráfica de Gantt semanal con actividades por alumno para Trabajo Terminal I.}
    \label{fig:proto-gantt-semanal-tt1}
\end{figure}

\subsection{Cronograma semanal de actividades de Trabajo Terminal II}

[Describa la planificación semanal de TT II, indicando la distribución de tareas entre los alumnos del equipo].

\begin{figure}[htbp]
    \centering
    % Inserte aquí su gráfica de Gantt exportada en PDF o PNG:
    % \includegraphics[width=0.95\textwidth]{figuras/proyecto/gantt-semanal-tt2.pdf}
    \fbox{\parbox[c][5cm][c]{0.9\textwidth}{\centering\textbf{[Espacio para Gráfica de Gantt Semanal de Trabajo Terminal II]}\\[0.5em]
    \small Inserte aquí el cronograma gráfico generado externamente mediante:\\
    \texttt{\textbackslash includegraphics[width=0.9\textbackslash textwidth]\{figuras/proyecto/gantt-semanal-tt2.pdf\}}}}
    \caption{Gráfica de Gantt semanal con actividades por alumno para Trabajo Terminal II.}
    \label{fig:proto-gantt-semanal-tt2}
\end{figure}

% ------------------------------------------------------------------------------
% CONFORMIDAD DEL CRONOGRAMA DE ACTIVIDADES (FIRMAS DE ALUMNOS Y ASESORES)
% ------------------------------------------------------------------------------
\vspace{1.5cm}
\begin{center}
    {\fontsize{11}{13}\bfseries\sffamily Conformidad del Cronograma de Actividades}\\[0.3em]
    {\fontsize{9}{11}\sffamily Con la firma de esta sección, los alumnos y asesores avalan el calendario de actividades y la viabilidad del proyecto.}\\[1.0em]
    \bloqueFirmasProtocolo[0.8cm]
\end{center}
